\documentclass[preprint,12pt,authoryear]{elsarticle}

\usepackage{amssymb}
\usepackage{amsthm}
\usepackage{lineno}

\usepackage[colorlinks=true,urlcolor=red,citecolor=blue%
anchorcolor=black,linkcolor=black]{hyperref}
\usepackage{comment}
\usepackage{soul}
\usepackage{color}
\usepackage{amsmath}
\usepackage{amsbsy}
\usepackage{dsfont}
\usepackage{subfigure}
\usepackage{threeparttable}
\usepackage{multirow}
\newcommand{\tabincell}[2]{\begin{tabular}{@{}#1@{}}#2\end{tabular}}
\usepackage{colortbl, booktabs}
\usepackage[dvipsnames]{xcolor}
\usepackage{algorithm}
\usepackage{algorithmic}
\usepackage{eurosym}
\usepackage{fullpage}
\usepackage{listings}
\usepackage{graphicx}
\usepackage{tikz}
\usetikzlibrary{shapes.geometric, arrows.meta, positioning, fit, backgrounds}

%% 中文支持（使用 XeLaTeX 或 LuaLaTeX 编译）
\usepackage{ctex}

% Listings style for Python code
\lstset{
  language=Python,
  basicstyle=\ttfamily\scriptsize,
  keywordstyle=\color{blue},
  commentstyle=\color{OliveGreen},
  stringstyle=\color{red},
  numbers=left,
  numberstyle=\tiny\color{gray},
  frame=single,
  breaklines=true,
  captionpos=b
}

% Revision commands
\newcommand{\lwnotes}[1]{\textcolor{cyan}{(Note: #1)}}
\newcommand{\lwedits}{\textcolor{cyan}}

\journal{IEEE Access}

\begin{document}

\begin{frontmatter}

\title{基于多Agent协作的自动化文献综述系统：\\聚焦底层原理分析与反幻觉保障}

\author[label1]{LeiWang Zhang}

\begin{abstract}
本文设计并实现了一套基于多Agent协作的自动化文献综述系统，覆盖从文献检索到结构化综述报告生成的全流程。系统采用协调者-代理（Coordinator-Agent）架构，由四个专业化Agent——检索（Search）、筛选（Screen）、聚类（Cluster）、摘要（Summary）——组成四阶段流水线，以层级研究问题（RQ）树驱动整个流程。核心创新包括：（1）三级筛选机制，结合基于嵌入的余弦相似度与LLM辅助的边界区域判定；（2）GPU加速的PCA降维与HDBSCAN聚类实现语义主题发现；（3）两阶段LLM报告生成策略，通过反幻觉提示词工程强制要求公式级分析、推导链条和第一性原理理论极限推导；（4）层级RQ树驱动筛选-聚类-摘要全流程。系统集成五个免费学术数据库（arXiv、PubMed、DBLP、Europe PMC、OpenAlex），支持命令行和Web界面两种交互方式。实验结果表明，系统能在数分钟内处理数百篇文献，生成包含强制公式解读和定量对比的结构化综述报告。
\end{abstract}

\begin{keyword}
多Agent系统 \sep 文献综述自动化 \sep 大语言模型 \sep 语义聚类 \sep 反幻觉 \sep 提示词工程
\end{keyword}

\end{frontmatter}

\section{引言}
\label{sec:intro}

科学出版物的快速增长给研究人员开展系统性文献综述带来了巨大挑战。据估计，每年各学科发表的新学术论文超过500万篇 \citep{bornmann-etal-2021}。一篇系统性文献综述通常要求研究人员手动筛选数千篇论文、识别相关主题并综合研究结果——这一过程可能耗时数周甚至数月。这种劳动密集型过程不仅耗时，还容易受到选择偏见和覆盖不全的影响。

近年来，大语言模型（LLM）和自然语言处理技术的进展为自动化文献综述的各个阶段开辟了新的可能。然而，现有方法往往只产生浅层摘要，仅仅罗列方法名称，而不解释底层原理、数学公式或性能差异背后的物理原因。一个真正有用的自动化综述应该超越表面描述，提供公式级分析、推导链条和第一性原理推理。

本文提出了一套多Agent自动化文献综述系统，通过以下四项关键创新解决上述挑战：

\begin{enumerate}
    \item \textbf{协调者-代理架构}：由四个专业化Agent（检索、筛选、聚类、摘要）组成顺序流水线，以层级研究问题（RQ）树驱动，通过检查点恢复机制保证数据流的连贯性和质量控制。
    \item \textbf{三级筛选机制}：结合基于嵌入的余弦相似度与LLM辅助的边界区域判定，在相关性的吞吐量和精度之间取得平衡。
    \item \textbf{GPU加速的语义聚类}：采用PCA降维和HDBSCAN聚类，支持自适应参数调节和领域一致性过滤，生成有意义的主题簇。
    \item \textbf{聚焦原理的报告生成}：通过反幻觉提示词工程强制要求公式解读、推导链条、理论极限分析和严格引用验证，生成分析"方法为何有效"而非仅仅罗列"方法是什么"的综述报告。
\end{enumerate}

系统集成五个免费学术数据库——arXiv、PubMed、DBLP、Europe PMC和OpenAlex——基本操作无需API密钥。支持命令行和Web界面两种交互方式，便于各学科研究人员使用。

本文其余部分组织如下：第~\ref{sec:lit-rev}~节综述相关工作；第~\ref{sec:model}~节描述系统架构和形式化问题定义；第~\ref{sec:method}~节详述各Agent的技术实现；第~\ref{sec:exp}~节展示实验结果和分析；第~\ref{sec:concl}~节总结全文。

\section{文献综述}
\label{sec:lit-rev}

本节从三个方面综述相关工作：自动化文献综述系统、研究任务的多Agent架构、以及LLM文本生成中的反幻觉技术。

\subsection{自动化文献综述系统}

文献综述过程的自动化日益受到研究关注。早期方法聚焦于基于关键词的检索和文献计量分析 \citep{white-mccain-1998}。近年来，\citet{wang-etal-2020}提出了使用抽象摘要的自动化综述生成框架。\citet{alaofi-etal-2023}探索了使用大语言模型生成文献综述的方法，展示了LLM生成综述的潜力和风险。

\citet{kostousov-2025}在IEEE Access上提出了AI增强的文献综述方法论，结合语义聚类和摘要生成。我们的系统在此方法论基础上引入了多项关键增强：（1）层级RQ驱动的流水线替代扁平主题聚类；（2）GPU加速的嵌入计算；（3）聚焦底层原理的提示词工程，强制要求公式级分析。

\subsection{多Agent架构}

多Agent系统已广泛应用于复杂任务分解。\citet{park-etal-2023}展示了基于LLM的Agent在协作场景中的有效性。\citet{hong-etal-2024}提出了MetaGPT，一个用于多Agent协作的元编程框架。本系统采用的协调者-代理模式受到这些工作启发，但专门针对文献综述领域进行了优化，引入了领域特定的质量门禁和检查点机制。

\subsection{LLM生成中的反幻觉技术}

幻觉——生成事实上不正确或虚构的内容——仍然是LLM应用中的关键挑战 \citep{ji-etal-2023}。在文献综述场景中，幻觉主要表现为虚构引用、错误的方法描述和缺乏依据的论断。\citet{shuster-etal-2021}证明检索增强生成（RAG）可以通过将生成内容锚定在检索到的文档上来减少幻觉。本系统采用多层反幻觉策略：严格的引用格式强制（附正确/错误示例）、禁止空泛表述、强制公式级推理，以及回退报告机制实现优雅降级。

\section{系统架构}
\label{sec:model}

\subsection{系统概述}

系统采用协调者-代理（Coordinator-Agent）架构，由中心化的\texttt{Coordinator}编排四个专业化Agent按顺序执行流水线。数据持久化和Agent间通信通过\texttt{SharedWorkspace}（共享工作区）管理，它为聚类和摘要提供按主题隔离的存储，同时全局共享文献和嵌入向量。

图~\ref{fig:architecture}展示了系统的整体架构。

\begin{figure}[htbp]
    \centering
    \begin{tikzpicture}[
        box/.style={rectangle, draw, rounded corners, minimum width=2.6cm, minimum height=1cm, align=center, font=\small},
        agent/.style={rectangle, draw, rounded corners, fill=blue!10, minimum width=2.2cm, minimum height=1cm, align=center, font=\small},
        arr/.style={-{Stealth[length=2.5mm]}, thick},
        darr/.style={-{Stealth[length=2.5mm]}, thick, dashed}
    ]
        % 用户输入
        \node[box, fill=yellow!20] (input) at (3,5) {用户输入\\研究主题};

        % 协调者
        \node[box, fill=red!10] (coord) at (3,3.5) {Coordinator\\流水线编排};

        % RQ管理器
        \node[box, fill=orange!10] (rq) at (8,3.5) {RQ Manager\\层级研究问题};

        % 四个Agent（均匀排列）
        \node[agent] (search)  at (0,1.5) {Search\\检索Agent};
        \node[agent] (screen)  at (2.8,1.5) {Screen\\筛选Agent};
        \node[agent] (cluster) at (5.6,1.5) {Cluster\\聚类Agent};
        \node[agent] (summary) at (8.4,1.5) {Summary\\摘要Agent};

        % 工作区
        \node[rectangle, draw, rounded corners, fill=green!10, minimum width=10cm, minimum height=0.8cm, align=center, font=\small] (ws) at (4.2,0) {Shared Workspace：文献 $\cdot$ 嵌入 $\cdot$ 聚类 $\cdot$ 报告};

        % 箭头：用户→协调者
        \draw[arr] (input) -- (coord);

        % 箭头：协调者→第一个Agent（链式，不扇形展开避免穿模）
        \draw[arr] (coord.south) -- ++(0,-0.4) -| (search.north);

        % 箭头：Agent顺序传递
        \draw[arr] (search) -- (screen);
        \draw[arr] (screen) -- (cluster);
        \draw[arr] (cluster) -- (summary);

        % 箭头：协调者→RQ管理器
        \draw[darr] (coord) -- (rq);

        % 箭头：各Agent→工作区
        \draw[arr] (search.south)  -- (ws.north -| search);
        \draw[arr] (screen.south)  -- (ws.north -| screen);
        \draw[arr] (cluster.south) -- (ws.north -| cluster);
        \draw[arr] (summary.south) -- (ws.north -| summary);

        % 箭头：RQ管理器→工作区（绕过summary，走右侧外边）
        \draw[darr] (rq.east) -- ++(0.4,0) |- ([xshift=0.5cm]ws.north east) -- (ws.north east);
    \end{tikzpicture}
    \caption{多Agent文献综述系统架构。Coordinator链式编排四个Agent，RQ Manager提供层级研究问题驱动筛选、聚类和摘要阶段。}
    \label{fig:architecture}
\end{figure}

\subsection{流水线阶段}

处理流水线由四个顺序阶段组成，每个阶段由专门的Agent负责：

\begin{table}[htbp]
    \centering
    \caption{流水线阶段及其功能}
    \label{tbl:pipeline}
    \begin{tabular}{llll}
    \toprule
    阶段 & Agent & 核心功能 & 关键技术 \\
    \midrule
    检索 & SearchAgent  & 多数据库文献检索 & arXiv, PubMed, DBLP, OpenAlex \\
    筛选 & ScreenAgent  & 相关性过滤 & BGE嵌入, 余弦相似度 \\
    聚类 & ClusterAgent & 主题发现 & GPU-PCA, HDBSCAN \\
    摘要 & SummaryAgent & 报告生成 & 两阶段LLM调用 \\
    \bottomrule
    \end{tabular}
\end{table}

\subsection{层级研究问题}

本系统的一项关键设计是使用层级研究问题（RQ）树驱动整个流水线。RQ树包含三个层级：

\begin{itemize}
    \item \textbf{第一级（宏观维度）}：方法论、应用、趋势等大方向，定义综述的整体范围。
    \item \textbf{第二级（中观分析）}：每个维度下的具体分析角度，如方法类型、应用领域、演进时间线。
    \item \textbf{第三级（微观深入）}：技术细节，包括核心数学公式与推导链条、第一性原理理论极限、模型假设与有效性边界。
\end{itemize}

RQ树根据研究主题自动生成，以JSON文件持久化。它为检索阶段提供上下文查询、为筛选阶段提供相关性标准、为聚类阶段提供语义锚点、为摘要阶段提供结构化Prompt。

\subsection{BaseAgent抽象}

所有Agent继承自通用的\texttt{BaseAgent}抽象基类，该基类提供：
\begin{itemize}
    \item LLM调用管理，通过全局\texttt{asyncio.Semaphore(5)}进行速率限制
    \item 标准化结果格式（\texttt{AgentResult}），包含成功/错误/指标字段
    \item 工作区集成，支持数据存储和检索
    \item 结构化日志，支持进度追踪
    \item 模板方法模式：子类实现\texttt{execute()}方法，基类负责编排
\end{itemize}

\subsection{共享工作区}

\texttt{SharedWorkspace}提供两级隔离的持久化存储：
\begin{itemize}
    \item \textbf{全局共享}：文献记录和嵌入向量在所有研究主题间共享，支持复用。
    \item \textbf{主题隔离}：聚类、摘要和报告按研究主题隔离。
    \item \textbf{检查点支持}：工作区支持创建和恢复检查点，实现工作流中断和恢复。
\end{itemize}

\section{技术实现}
\label{sec:method}

本节详述各流水线阶段的技术实现。

\subsection{多源异构文献检索}
\label{sec:method-search}

\texttt{SearchAgent}通过以下步骤实现多数据库文献检索：

\textbf{查询构建。}给定研究主题，Agent首先调用LLM生成结构化搜索策略。提示词强制要求简洁的关键词组合（每条查询3--5个词）、数据库特定语法，以及通过领域限定词消除歧义。若LLM未能返回有效查询，回退策略将从研究主题中生成关键词两两和三三组合。

\textbf{多数据库搜索。}系统集成五个学术数据库，均无需API密钥即可访问：

\begin{table}[htbp]
    \centering
    \caption{支持的学术数据库}
    \label{tbl:databases}
    \begin{tabular}{llcc}
    \toprule
    数据库 & 领域 & 摘要 & 需要API Key \\
    \midrule
    arXiv       & 物理、计算机、数学     & 有 & 否 \\
    PubMed      & 医学、生物学           & 有 & 否 \\
    DBLP        & 计算机科学             & 无 & 否 \\
    Europe PMC  & 生命科学、生物医学     & 有 & 否 \\
    OpenAlex    & 跨学科                 & 无 & 否 \\
    \bottomrule
    \end{tabular}
\end{table}

每个数据库实施独立的速率限制，其中arXiv要求查询间3秒延迟以避免被封禁。若某次arXiv查询返回零结果，则跳过后续arXiv查询。

\textbf{去重与标准化。}检索到的论文基于标题标准化比较进行去重。每篇论文标准化为\texttt{LiteratureRecord}数据结构，包含标题、作者、摘要、年份、来源、DOI、期刊/会议和引用数字段。

\subsection{三级相关性筛选}
\label{sec:method-screen}

\texttt{ScreenAgent}实现三级筛选机制，在吞吐量和精度之间取得平衡。

\textbf{第一级：基于嵌入的余弦相似度。}
对每篇论文，计算其嵌入（标题+摘要）与每个RQ嵌入（问题+关键词）之间的余弦相似度。设$\mathbf{p}_i \in \mathbb{R}^d$为论文$i$的嵌入，$\mathbf{r}_j \in \mathbb{R}^d$为RQ $j$的嵌入，论文$i$的相关性分数为：
\begin{equation}
    s_i = \max_{j} \frac{\mathbf{p}_i \cdot \mathbf{r}_j}{\|\mathbf{p}_i\| \cdot \|\mathbf{r}_j\|}
    \label{eq:cosine-similarity}
\end{equation}
通过矩阵乘法$\mathbf{P} \mathbf{R}^\top$一次性计算，其中$\mathbf{P} \in \mathbb{R}^{n \times d}$和$\mathbf{R} \in \mathbb{R}^{m \times d}$分别为论文和RQ的嵌入矩阵。

\textbf{第二级：基于阈值的分类。}
论文根据相似度分数分为三个区域：
\begin{itemize}
    \item \textbf{自动通过}（$s_i \geq \theta_{\max}$）：直接标记为相关
    \item \textbf{自动拒绝}（$s_i < \theta_{\min}$）：直接标记为不相关
    \item \textbf{边界区域}（$\theta_{\min} \leq s_i < \theta_{\max}$）：送入LLM判定
\end{itemize}
其中$\theta_{\min} = 0.68$、$\theta_{\max} = 0.80$为可配置阈值。

\textbf{第三级：LLM辅助边界判定。}
对边界区域论文，LLM从三个维度进行加权评分：
\begin{equation}
    w_i = 0.5 \cdot \text{topic}_i + 0.3 \cdot \text{method}_i + 0.2 \cdot \text{timeliness}_i
    \label{eq:weighted-score}
\end{equation}
当且仅当$w_i \geq 3.0$时，论文被标记为相关。边界论文按相似度分数排序，限制最多150次LLM调用以控制成本。

\textbf{嵌入模型选择。}系统自动检测论文内容的语言，选择合适的嵌入模型：
\begin{itemize}
    \item 中文文本：BGE-small-zh-v1.5（768维，本地GPU推理）
    \item 英文文本：all-MiniLM-L6-v2（384维，本地GPU推理）
    \item 回退方案：远程GLM embedding-3 API（1024维）
\end{itemize}

\subsection{GPU加速语义聚类}
\label{sec:method-cluster}

\texttt{ClusterAgent}对相关论文进行多步骤语义聚类。

\textbf{降维。}
采用两步降维策略。给定嵌入矩阵$\mathbf{E} \in \mathbb{R}^{n \times d}$：
\begin{enumerate}
    \item \textbf{第一步：}PCA将维度从$d$降至50，GPU可用时使用PyTorch CUDA SVD分解：
    \begin{equation}
        \mathbf{E}_{50} = \mathbf{U}_{(:, 1:50)} \cdot \text{diag}(\mathbf{S}_{1:50})
        \label{eq:pca-reduction}
    \end{equation}
    其中$\mathbf{E} - \bar{\mathbf{E}} = \mathbf{U} \mathbf{S} \mathbf{V}^\top$为SVD分解。
    \item \textbf{第二步：}t-SNE从50维降至2维用于可视化，困惑度根据样本量自适应：
    \begin{equation}
        \text{perplexity} = \min(30, \max(5, n / 10))
        \label{eq:perplexity}
    \end{equation}
\end{enumerate}
对于小数据集（$n < 10$），直接使用PCA降至2维，跳过t-SNE。

\textbf{HDBSCAN聚类。}
HDBSCAN算法在降维空间中聚类论文，自适应参数：
\begin{equation}
    \text{min\_cluster\_size} = \max(2, n / 10)
    \label{eq:min-cluster}
\end{equation}
无法分配到任何簇的论文标记为噪声（$-1$）。当HDBSCAN不可用时，KMeans作为回退方案。

\textbf{聚类标签生成。}
对每个聚类，LLM生成描述性标签、核心主题、代表性论文、子主题，以及该簇论文共享的核心公式或数学框架和底层物理/数学原理。

\textbf{领域一致性过滤。}
通过比较聚类论文标题与研究主题关键词的重叠率进行过滤。主题关键词匹配率低于30\%的聚类被移除。

\subsection{聚焦原理的报告生成}
\label{sec:method-summary}

\texttt{SummaryAgent}通过精心设计的提示词工程流水线生成结构化综述报告。

\textbf{两阶段生成。}
最终报告通过两次LLM调用生成，以绕过输出token限制：
\begin{itemize}
    \item \textbf{第一部分：}引言 + 核心技术原理（按聚类分析）
    \item \textbf{第二部分：}应用分析 + 技术演进 + 理论极限 + 突破路径 + 结论
\end{itemize}

\textbf{反幻觉提示词设计。}
提示词在多个层级强制执行严格规则：

\begin{table}[htbp]
    \centering
    \caption{反幻觉提示词工程层级}
    \label{tbl:antihallucination}
    \begin{tabular}{lp{7cm}}
    \toprule
    层级 & 强制规则 \\
    \midrule
    引用格式 & 强制格式：[作者, 年份, 标题]。提供错误示例。禁止引用列表外论文。 \\
    反空泛表述 & "取得了较好效果"若无具体指标值则禁止使用。每个论断必须引用定量数据。 \\
    公式深度 & 每个方法必须列出核心公式，逐符号解释（含义、量纲、取值范围）。必须说明推导来源。 \\
    原理级对比 & 禁止"A优于B"而不解释公式/物理层面的原因。 \\
    理论极限 & 必须从第一性原理（物理定律、数学定理）推导性能上下界。 \\
    \bottomrule
    \end{tabular}
\end{table}

\textbf{回退机制。}
始终生成基于模板的回退报告作为安全网。若LLM失败（如3次重试后仍遇到速率限制），回退报告作为最终输出。

\textbf{语言一致性。}
系统检测研究主题的语言，在整个报告中强制使用一致的语言。中文主题生成中文报告，英文主题生成英文报告。

\section{系统演示与性能分析}
\label{sec:exp}

本节展示系统的性能特征和能力演示。

\subsection{实验环境}

系统基于Python 3.10+实现，主要依赖如下：
\begin{itemize}
    \item \textbf{LLM框架：}LangChain + OpenAI兼容接口
    \item \textbf{默认LLM：}智谱AI GLM-4-Plus
    \item \textbf{嵌入模型：}BGE-small-zh-v1.5（本地）、all-MiniLM-L6-v2（本地）、GLM embedding-3（远程）
    \item \textbf{降维：}PyTorch CUDA PCA、scikit-learn t-SNE
    \item \textbf{聚类：}HDBSCAN
    \item \textbf{向量索引：}FAISS（GPU $>$ CPU $>$ NumPy回退链）
    \item \textbf{异步框架：}asyncio + aiohttp
    \item \textbf{Web界面：}Gradio
\end{itemize}

\subsection{性能优化}

表~\ref{tbl:optimizations}总结了系统中实现的关键性能优化。

\begin{table}[htbp]
    \centering
    \caption{性能优化技术}
    \label{tbl:optimizations}
    \begin{tabular}{lp{7cm}}
    \toprule
    优化项 & 实现方式 \\
    \midrule
    LLM并发控制 & 全局\texttt{asyncio.Semaphore(5)}，所有LLM调用通过基类统一限流 \\
    GPU加速 & PyTorch CUDA PCA降维，FAISS向量索引（优先GPU） \\
    批量处理 & 相似度矩阵一次矩阵乘法完成；批量嵌入生成 \\
    LLM客户端缓存 & 按模型+参数组合缓存ChatOpenAI实例 \\
    嵌入复用 & 聚类阶段复用已缓存的嵌入，仅计算缺失部分 \\
    异步并行 & 通过\texttt{asyncio.gather}并行生成聚类摘要 \\
    \bottomrule
    \end{tabular}
\end{table}

\subsection{流水线执行流程}

对于典型研究主题，执行流程如下：

\begin{enumerate}
    \item \textbf{RQ生成}（$<1$秒）：RQ管理器根据研究主题生成三级RQ树。以"降低晶体管亚阈值摆幅的方法"为主题，RQ树将包含RQ1（方法论）、RQ2（应用）、RQ3（趋势），每个维度下含聚焦核心公式、推导链条、理论极限和模型假设的第二级和第三级子问题。

    \item \textbf{检索}（2--5分钟）：检索Agent生成8--10条查询，在多个数据库中搜索。考虑速率限制（arXiv 3秒延迟，其他0.5秒），一次典型运行检索200--500篇论文。

    \item \textbf{筛选}（1--3分钟）：嵌入计算是此阶段的主要耗时。余弦相似度矩阵通过单次矩阵运算完成。边界区域论文（通常占候选论文的20--40\%）送入LLM判定。

    \item \textbf{聚类}（30秒--2分钟）：PCA和t-SNE由GPU加速。HDBSCAN在CPU上运行。聚类标签生成涉及并行LLM调用。

    \item \textbf{摘要}（2--5分钟）：两次顺序LLM调用生成最终报告。聚类摘要并行生成。始终生成回退报告作为备份。
\end{enumerate}

\subsection{输出质量}

生成的报告遵循结构化格式，包含以下章节：
\begin{enumerate}
    \item \textbf{引言：}研究背景、从物理/数学视角阐述核心挑战
    \item \textbf{核心技术原理：}按聚类分析，包含：
    \begin{itemize}
        \item 核心公式及逐符号解读
        \item 从基本定律出发的推导链条
        \item 关键参数和理论极限
        \item 定量比较表格
    \end{itemize}
    \item \textbf{实际技术效果：}按应用场景展示公式适配的量化指标
    \item \textbf{技术演进与物理极限：}各代方法的公式演进、从第一性原理推导理论极限
    \item \textbf{根因分析与突破路径：}公式层面的瓶颈识别、未解决的基础性问题
    \item \textbf{结论：}核心公式级发现和定量总结
\end{enumerate}

\subsection{Web界面}

系统提供基于Gradio的Web界面，包含五个功能Tab：

\begin{figure}[htbp]
    \centering
    \begin{tikzpicture}[
        tab/.style={rectangle, draw, rounded corners, fill=blue!8, minimum width=2cm, minimum height=1.4cm, align=center, font=\scriptsize}
    ]
        \node[tab] (t1) at (0,0)    {新建综述\\输入主题\\执行流水线};
        \node[tab] (t2) at (2.6,0)  {文献库\\浏览和筛选\\论文列表};
        \node[tab] (t3) at (5.2,0)  {聚类可视化\\2D散点图\\主题分布};
        \node[tab] (t4) at (7.8,0)  {综述报告\\查看生成的\\报告文档};
        \node[tab] (t5) at (10.4,0) {系统状态\\工作区统计\\RQ树};
    \end{tikzpicture}
    \caption{Web界面Tab结构}
    \label{fig:webui}
\end{figure}

界面支持实时流水线执行与进度日志、多主题管理，以及聚类结果的数据可视化。

\section{结论}
\label{sec:concl}

本文提出了一套多Agent自动化文献综述系统，解决了三个关键挑战：（1）手动文献筛选的可扩展性问题；（2）自动化分析的深度问题；（3）LLM生成内容的可靠性问题。

系统的核心创新包括：层级RQ驱动的流水线确保各阶段数据流的连贯性；三级筛选机制在吞吐量和精度之间取得平衡；GPU加速的语义聚类支持自适应参数调节；聚焦原理的报告生成策略通过多层反幻觉保障确保输出质量。

本系统的一个显著特点是强调公式级分析。与现有自动化综述工具产生浅层摘要不同，本系统通过精心设计的提示词强制要求每个方法必须附带核心公式及逐符号解释、从基本定律出发的推导链条、以及从第一性原理推导的理论性能极限。这种分析深度得益于层级RQ结构中包含的专门针对数学公式、物理模型和假设有效性的微观级问题。

系统集成五个免费学术数据库无需API密钥，支持命令行和Web界面两种交互方式，并通过检查点机制实现工作流中断恢复。GPU加速PCA、FAISS向量索引和并行LLM调用等性能优化使系统能在数分钟内处理数百篇论文。

未来工作将聚焦于：（1）引入全文PDF分析以克服当前仅依赖摘要的局限性，实现更深入的公式和推导提取；（2）扩展RQ树以支持不同研究领域的领域特定模板；（3）实现迭代优化机制，系统自动识别覆盖缺口并扩展检索范围；（4）开发自动化综述质量评估指标，包括公式准确性验证。



\section*{参考文献}
\bibliographystyle{elsarticle-harv}
\bibliography{bib-template}

\end{document}
